\documentclass{article}
\usepackage{amsmath}
\usepackage{amssymb}
\usepackage[margin=1in]{geometry}

\begin{document}

\begin{center}
{\LARGE CSE400 -- Fundamentals of Probability in Computing}

\vspace{0.5cm}

{\Large Lecture 11: Transformation of Random Variables}

\vspace{0.5cm}

{\large Lecture Scribe for Exam Revision}
\end{center}

\vspace{1cm}

\section*{Transformation of Random Variables}

\subsection*{Lecture Objective}

The lecture focuses on three connected goals:

\begin{itemize}
\item \textbf{Transformation of Random Variables}

Learn techniques to determine the distribution of a new random variable obtained from another.

\item \textbf{Function of Two Random Variables}

Study transformations when the new variable depends on two random variables.

\item \textbf{Illustrative Example}

Work through the case

\[
Z = X + Y
\]
\end{itemize}

These topics form the logical progression of the lecture:

single variable transformation $\rightarrow$ two variable transformation $\rightarrow$ example derivation.

\section*{1. Transformation of a Random Variable}

\subsection*{Basic Setup}

We start with a known random variable:

\[
X \sim f_X(x)
\]

Assumption:

The PDF of $X$ is known a priori.

We define a new random variable

\[
Y = g(X)
\]

Goal:

Find the PDF of $Y$:

\[
f_Y(y)
\]

This means we must convert the probability description of $X$ into the probability description of $Y$.

\subsection*{Case 1: Monotonically Increasing Transformation}

Assume:

\[
Y = g(X)
\]

and $g(x)$ is monotonically increasing.

This assumption is important because it guarantees:

\begin{itemize}
\item the inverse function exists
\item each $y$ corresponds to exactly one $x$
\end{itemize}

Thus

\[
x = g^{-1}(y)
\]

\subsection*{Step 1 — Express the CDF of $Y$}

Start from the definition of CDF:

\[
F_Y(y) = P(Y \le y)
\]

Substitute $Y = g(X)$:

\[
F_Y(y) = P(g(X) \le y)
\]

Since $g$ is increasing, the inequality direction does not change, therefore

\[
P(g(X) \le y) = P(X \le g^{-1}(y))
\]

Thus

\[
F_Y(y) = F_X(g^{-1}(y))
\]

This step converts the probability into a form involving the known distribution of $X$.

\subsection*{Step 2 — Differentiate to obtain the PDF}

The PDF is the derivative of the CDF:

\[
f_Y(y) = \frac{d}{dy}F_Y(y)
\]

Substitute the CDF expression:

\[
f_Y(y) = \frac{d}{dy}[F_X(g^{-1}(y))]
\]

Apply the chain rule:

\[
f_Y(y) = f_X(g^{-1}(y)) \cdot \frac{d}{dy}g^{-1}(y)
\]

Since

\[
x = g^{-1}(y)
\]

we can write

\[
f_Y(y) = f_X(x)\frac{dx}{dy}
\]

evaluated at

\[
x = g^{-1}(y)
\]

\subsection*{Case 2: Monotonically Decreasing Transformation}

Now assume $g(x)$ is decreasing.

The key difference is the inequality direction flips.

\subsection*{Step 1 — CDF expression}

Start again from

\[
F_Y(y) = P(Y \le y)
\]

Substitute transformation:

\[
F_Y(y) = P(g(X) \le y)
\]

Because the function is decreasing:

\[
g(X) \le y \Rightarrow X \ge g^{-1}(y)
\]

Thus

\[
F_Y(y) = P(X \ge g^{-1}(y))
\]

Using the complement rule:

\[
F_Y(y) = 1 - F_X(g^{-1}(y))
\]

\subsection*{Step 2 — Differentiate}

Differentiate the CDF:

\[
f_Y(y) = \frac{d}{dy}[1 - F_X(g^{-1}(y))]
\]

\[
f_Y(y) = -f_X(g^{-1}(y))\frac{d}{dy}g^{-1}(y)
\]

The negative sign appears because the transformation is decreasing.

\subsection*{Final General Formula}

To combine both cases, the PDF is written using an absolute value:

\[
f_Y(y) = f_X(x)\left|\frac{dx}{dy}\right|
\]

evaluated at

\[
x = g^{-1}(y)
\]

This formula ensures the PDF is always positive.

\section*{Example: Transformation of a Uniform Random Variable}

Given:

\[
X \sim U(-1,1)
\]

So the PDF is

\[
f_X(x) =
\begin{cases}
\frac{1}{2}, & -1 < x < 1 \\
0, & \text{otherwise}
\end{cases}
\]

Define the transformation:

\[
Y = g(X) = \sin\left(\frac{\pi X}{2}\right)
\]

Goal:

Find

\[
f_Y(y)
\]

\subsection*{Step 1 — Find the inverse transformation}

\[
y = \sin\left(\frac{\pi x}{2}\right)
\]

Solve for $x$:

\[
x = \frac{2}{\pi}\sin^{-1}(y)
\]

\subsection*{Step 2 — Compute derivative}

\[
\frac{dx}{dy} = \frac{2}{\pi}\frac{1}{\sqrt{1-y^2}}
\]

\subsection*{Step 3 — Apply transformation formula}

\[
f_Y(y) = f_X(x)\left|\frac{dx}{dy}\right|
\]

Substitute $f_X(x) = \frac{1}{2}$:

\[
f_Y(y) = \frac{1}{2}\times \frac{2}{\pi}\frac{1}{\sqrt{1-y^2}}
\]

Simplify:

\[
f_Y(y) = \frac{1}{\pi\sqrt{1-y^2}}
\]

\subsection*{Step 4 — Determine support}

From the transformation:

if $x = -1$

\[
y = \sin(-\pi/2) = -1
\]

if $x = 1$

\[
y = \sin(\pi/2) = 1
\]

Thus

\[
-1 < y < 1
\]

Final PDF:

\[
f_Y(y) =
\begin{cases}
\frac{1}{\pi\sqrt{1-y^2}}, & -1 < y < 1 \\
0, & \text{otherwise}
\end{cases}
\]

\section*{Function of Two Random Variables}

The lecture next considers the transformation

\[
Z = X + Y
\]

Goal:

Find the PDF

\[
f_Z(z)
\]

\section*{Deriving the Distribution of $Z = X + Y$}

Start with the CDF definition:

\[
F_Z(z) = P(Z \le z)
\]

Substitute the definition of $Z$:

\[
F_Z(z) = P(X + Y \le z)
\]

This probability corresponds to a region in the $x$-$y$ plane.

To compute it we integrate the joint PDF over that region.

\subsection*{Horizontal strip integration}

\[
F_Z(z)=\int_{y=-\infty}^{\infty}\int_{x=-\infty}^{z-y}f_{XY}(x,y)\,dx\,dy
\]

inner limit: $x \le z-y$ because $x+y \le z$

\subsection*{Vertical strip integration}

\[
F_Z(z)=\int_{x=-\infty}^{\infty}\int_{y=-\infty}^{z-x}f_{XY}(x,y)\,dy\,dx
\]

Both integrals describe the same probability region.

\section*{Special Cases Mentioned}

Find $f_Z(z)$ in general.

If $X$ and $Y$ are independent.

If

\[
X \sim N(0,1), \quad Y \sim N(0,1)
\]

then

\[
Z = X + Y
\]

and it can be shown

\[
Z \sim N(0,2)
\]

If $X$ and $Y$ are exponential random variables with parameter $\lambda$, compute $f_Z(z)$.

\section*{Key Logical Flow of the Lecture}

\begin{itemize}
\item Start with known distribution of $X$.
\item Define a new random variable $Y = g(X)$.
\item Express CDF of $Y$ using probabilities of $X$.
\item Differentiate to obtain the PDF of $Y$.
\item Handle increasing vs decreasing functions.
\item Apply the formula to a worked example.
\item Extend the concept to two random variables.
\item Use joint PDF integration to compute distributions like $Z = X + Y$.
\end{itemize}

\end{document}